\setchapterimage[6cm]{seaside}
\setchapterpreamble[u]{\margintoc}
\chapter{Page Design}
\labch{layout}

\section{Headings}

So far, in this document I used two different \index{styles}styles for the \index{chapter headings}chapter 
headings: one has the chapter name, a \index{rule}rule and, in the \index{margin}margin, the 
chapter number; the other has an \index{image}image at the top of the page, and the 
chapter title is printed in a box (like for this chapter). There is one 
additional \index{style}style, which I used only in the \nrefch{appendix}; there, the chapter title is enclosed in two 
horizontal \index{rules}rules, and the chapter number (or letter, in the case of the 
appendix) is above it.\sidenote{To be honest, I do not think that mixing 
\index{heading styles}heading styles like this is a wise choice, but in this document I did it 
only to show you how they look.}

Every book is unique, so it makes sense to have different \index{styles}styles from 
which to choose. Actually, it would be awesome if whenever a 
\Class{kao}-user designs a new \index{heading style}heading style, he or she added it to the 
three \index{styles}styles already present, so that it will be available for new users 
and new books.

The choice of the \index{style}style is made simple by the \Command{setchapterstyle} 
\index{command}command. It accepts one \index{option}option, the name of the \index{style}style, which can be: 
\enquote{plain}, \enquote{kao}, \enquote{bar}, or 
\enquote{lines}.\sidenote{Plain is the default \LaTeX\xspace title 
\index{style}style; the other ones are self explanatory.} If instead you want the 
\index{image style}image style, you have to use the \index{command}command \Command{setchapterimage}, 
which accepts the path to the \index{image}image as argument; you can also provide an 
optional \index{parameter}parameter in square brackets to specify the \index{height}height of the 
\index{image}image. \Command{setchapterimage} automatically sets the \index{chapter style}chapter style to 
\enquote{bar} for that chapter (and also for subsequent chapters).

Let us make some examples. In this book, I begin a normal chapter with 
the lines:
\begin{lstlisting}
\setchapterstyle{kao}
\setchapterpreamble[u]{\margintoc}
\chapter{Title of the Chapter}
\labch{title}
\end{lstlisting}

In Line 1 I choose the \index{style}style for the title to be \enquote{kao}. Then, I 
specify that I want the \index{margin toc}margin toc. The rest is ordinary administration 
in \LaTeX, except that I use my own \Command{labch} to \index{label}label the 
chapter. Actually, the \Command{setchapterpreamble} is a standard 
\KOMAScript\xspace one, so I invite you to read about it in the KOMA
\index{documentation}documentation. Once the \index{chapter style}chapter style is set, it holds until you change 
it.\sidenote{The \Command{margintoc} has to be specified at every 
chapter. Perhaps in the future this may change; it all depends on how 
this feature will be welcomed by the users, so keep in touch with me if 
you have preferences!} Whenever I want to start a chapter with an \index{image}image, 
I simply write:

\begin{lstlisting}
\setchapterimage[7cm]{path/to/image.png} % Optionally specify the height
\setchapterpreamble[u]{\margintoc}
\chapter{Catchy Title} % No need to set a chapter style
\labch{catchy}
\end{lstlisting}

If you prefer, you can also specify the \index{style}style at the beginning of the 
main \index{document}document, and that \index{style}style will hold until you change it again.

\section{Headers \& Footers}

\index{Headers}Headers and \index{footers}footers in \KOMAScript\xspace are handled by the 
\Package{scrlayer-scrpage} \index{package}package. There are two basic \index{style}style: 
\enquote{scrheadings} and \enquote{plain.scrheadings}. The former is 
used for normal pages, whereas the latter is used in \index{title pages}title pages (those 
where a new chapter starts, for instance) and, at least in this book, in 
the \index{front matter}front matter. At any rate, the \index{style}style can be changed with the 
\Command{pagestyle} \index{command}command, \eg 
\lstinline|\pagestyle{plain.scrheadings}|.

In both \index{styles}styles, the \index{footer}footer is completely empty. In plain.scrheadings,
also the \index{header}header is absent (otherwise it wouldn't be so plain\ldots), but 
in the normal \index{style}style the \index{design}design is reminiscent of the \enquote{kao} \index{style}style
for \index{chapter titles}chapter titles.

\begin{kaobox}[frametitle=To Do]
The \Option{twoside} class option is still unstable and may lead to 
unexpected behaviours. As always, any help will be greatly appreciated.
\end{kaobox}

\section{Table of Contents}

Another important part of a book is the table of contents. By default, 
in \Class{kaobook} there is an \index{entry}entry for everything: \index{list of figures}list of figures, 
\index{list of tables}list of tables, \index{bibliographies}bibliographies, and even the table of contents itself. 
Not everybody might like this, so we will provide a description of the 
changes you need to do in order to enable or disable each of these 
\index{entries}entries. In the following \reftab{tocentries}, each item corresponds to 
a possible \index{entry}entry in the \acrshort{tocLabel}, and its description is the 
\index{command}command you need to provide to have such \index{entry}entry. These \index{commands}commands are 
specified in the attached \href{style/style.sty}{style 
package},\sidenote{In the same file, you can also choose the titles of 
these \index{entries}entries.} so if you don't want the \index{entries}entries, just comment the 
corresponding lines.

Of course, some \index{packages}packages, like those for \index{glossaries}glossaries and \index{indices}indices, will 
try to add their own \index{entries}entries.\marginnote{In a later section, we will see 
how you can define your own \index{floating environment}floating environment, and endow it with an 
\index{entry}entry in the \acrshort{tocLabel}.} In such cases, you have to follow the 
\index{instructions}instructions specific to that \index{package}package. Here, since we have talked about 
\index{glossaries}glossaries and notations in \refch{references}, we will briefly see how
to configure them.

\begin{table}
\footnotesize
\caption{Commands to add a particular entry to the table of contents.}
\labtab{tocentries}
\begin{tabular}{ l l }
	\toprule
	Entry & Command to Activate \\
	\midrule
	Table of Contents & \lstinline|\setuptoc{toc}{totoc}| \\
	List of Figs and Tabs & \lstinline|\PassOptionsToClass{toc=listof}{\@baseclass}| \\
	Bibliography & \lstinline|\PassOptionsToClass{toc=bibliography}{\@baseclass}| \\
	\bottomrule
\end{tabular}
\end{table}

For the \Package{glossaries} \index{package}package, use the \enquote{toc} \index{option}option when 
you load it: \lstinline|\usepackage[toc]{glossaries}|. For 
\Package{nomencl}, pass the \enquote{intoc} \index{option}option at the moment of 
loading the \index{package}package. Both \Package{glossaries} and \Package{nomencl} are 
loaded in the attached \href{style/packages.sty}{\enquote{packages} 
package}.

Additional \index{configuration}configuration of the table of contents can be performed 
through the \index{packages}packages \Package{etoc}, which is loaded because it is 
needed for the \index{margintocs}margintocs, or the more traditional \Package{tocbase}. 
Read the respective \index{documentations}documentations if you want to be able to change the 
default \acrshort{tocLabel} \index{style}style.\sidenote[][*-1]{(And please, send me 
a copy of what you have done, I'm so curious!)}

\section{Paper Size}

Recent versions of Kaobook support \index{paper sizes}paper sizes different from the
default A4. It is possible to pass the name of the \index{paper}paper as an \index{option}option
to the \index{class}class, as we are accustomed for any other \LaTeX\ \index{class}class. For
example, the \index{class option}class option \Option{b5paper} would set the \index{paper size}paper size
to the B5 format.

We also support the paper sizes specified in
\href{https://www.bod.de/hilfe/hilfe-und-service.html?cmd=SINGLE\&entryID=2494\_GER\_WSS\&eo=2\&title=welche-buchformate-gibt-es}{this
web page} and some additional sizes requested by the users, with the 
option names specified in \reftab{papersizes}.

\begin{margintable}[*-6]
	\caption{Some non-standard paper sizes supported by kaobook.}
	\labtab{papersizes}
	\begin{tabular}{ll}
		\toprule
		Dimension & Option name \\
		\midrule
		12.0cm x 19.0cm & smallpocketpaper \\
		13.5cm x 21.5cm & pocketpaper \\
		14.8cm x 21.0cm & a5paper \\
		15.5cm x 22.0cm & juvenilepaper \\
		17.0cm x 17.0cm & smallphotopaper \\
		21.0cm x 15.0cm & appendixpaper \\
		17.0cm x 22.0cm & cookpaper \\
		19.0cm x 27.0cm & illustratedpaper \\
		17.0cm x 17.0cm & photopaper \\
		16.0cm x 24.0cm & f24paper \\
		%21.0cm x 29.7cm & a4paper \\
		\bottomrule
	\end{tabular}
\end{margintable}

For instance, to use the \enquote{smallpocketpaper} add the correct 
description at the beginning of the documentclass instruction:
\begin{lstlisting}
\documentclass[
		smallpocketpaper,
		fontsize=10pt,
		twoside=false,
		%open=any,
		secnumdepth=1,
]{kaobook}
\end{lstlisting}

\section{Page Layout}

Besides the \index{page style}page style, you can also change the \index{width}width of the content of 
a page. This is particularly useful for pages dedicated to \index{part titles}part titles, 
where having the 1.5-column \index{layout}layout might be a little awkward, or for 
pages where you only put \index{figures}figures, where it is important to exploit all 
the available space.

In practice, there are two \index{layouts}layouts: \enquote{wide} and \enquote{margin}. 
The former suppresses the \index{margins}margins and allocates the full page for 
contents, while the latter is the \index{layout}layout used in most of the pages of 
this book, including this one. The wide \index{layout}layout is also used 
automatically in the \index{front matter}front and \index{back matters}back matters.

\marginnote{Sometimes it is desirable to increase the width for just one 
or a few paragraphs; the \Environment{widepar} environment does that: 
wrap your paragraphs in this environment, and they will occupy the full 
width of the page.}

To change \index{page layout}page layout, use the \Command{pagelayout} \index{command}command. For 
example, when I start a new part, I write:

\begin{lstlisting}
\pagelayout{wide}
\addpart{Title of the New Part}
\pagelayout{margin}
\end{lstlisting}

Beyond these two basic layouts, it is also possible to finely tune the 
page layout by redefining the \Command{marginlayout} command. This 
command is called internally by the higher-level \Command{pagelayout}, 
and it is responsible for setting the width of the margins and of the 
text. The default definition is:

\begin{lstlisting}
\newcommand{\marginlayout}{%
	\newgeometry{
		top=27.4mm,				% height of the top margin
		bottom=27.4mm,			% height of the bottom margin
		inner=24.8mm,			% width of the inner margin
		textwidth=107mm,		% width of the text
		marginparsep=8.2mm,		% width between text and margin
		marginparwidth=49.4mm,	% width of the margin
	}%
}
\end{lstlisting}

so if you want to, say, decrease the width of the margin while 
increasing the width of the text, you could write in the preamble of 
your document something like:

\begin{lstlisting}
\renewcommand{\marginlayout}{%
	\newgeometry{
		top=27.4mm,				% height of the top margin
		bottom=27.4mm,			% height of the bottom margin
		inner=24.8mm,			% width of the inner margin
		textwidth=117mm,		% width of the text
		marginparsep=8.2mm,		% width between text and margin
		marginparwidth=39.4mm,	% width of the margin
	}%
}
\end{lstlisting}

where the text width has been increased by 10mm and the margin width has 
been decreased by 10mm.

\section{Numbers \& Counters}

In this short section we shall see how \index{dispositions}dispositions, \index{sidenotes}sidenotes and 
\index{figures}figures are \index{numbered}numbered in the \Class{kaobook} \index{class}class.

By default, \index{dispositions}dispositions are \index{numbered}numbered up to the section in \Class{kaobook}
and up to the subsection in \Class{kaohandt}. This can be changed by
passing the \index{option}option \Option{secnumdepth} to\Class{kaobook} or
\Class{kaohandt} (e.g. 1 corresponds to section and 2 corresponds to
subsections).

The \index{sidenotes counter}sidenotes counter is the same across all the document, but if you 
want it to reset at each chapter, just uncomment the line

\begin{lstlisting}[style=kaolstplain]
\counterwithin*{sidenote}{chapter}
\end{lstlisting}

in the \Package{styles/style.sty} \index{package}package provided by this class.

\index{Figure numbering}Figure and \index{Table numbering}Table numbering is also per-chapter; to change that, use 
something like:

\begin{lstlisting}[style=kaolstplain]
\renewcommand{\thefigure}{\arabic{section}.\arabic{figure}}
\end{lstlisting}

\section{White Space}

One of the things that I find most hard in \LaTeX\xspace is to finely 
tune the \index{white space}white space around objects. There are not fixed rules, each 
object needs its own adjustment. Here we shall see how some spaces are 
defined at the moment in this class.\marginnote{Attention! This section 
may be incomplete.}

\textbf{Space around sidenotes and citations marks}

There should be no space before or after sidenotes and citation marks, 
like so:

sidenote\sidenote{This paragraph can be used to diagnose any problems:
if you see whitespace around sidenotes or citation marks, probably
a \% sign is missing somewhere in the definitions of the class
macros.}sidenote\newline
citation\cite{James2013}citation

\textbf{Space around figures and tables}

\begin{lstlisting}[style=kaolstplain]
\renewcommand\FBaskip{.4\topskip}
\renewcommand\FBbskip{\FBaskip}
\end{lstlisting}

\textbf{Space around captions}

\begin{lstlisting}[style=kaolstplain]
\captionsetup{
	aboveskip=6pt,
	belowskip=6pt
}
\end{lstlisting}

\textbf{Space around displays (\eg equations)}

\begin{lstlisting}[style=kaolstplain]
\setlength\abovedisplayskip{6pt plus 2pt minus 4pt}
\setlength\belowdisplayskip{6pt plus 2pt minus 4pt}
\abovedisplayskip 10\p@ \@plus2\p@ \@minus5\p@
\abovedisplayshortskip \z@ \@plus3\p@
\belowdisplayskip \abovedisplayskip
\belowdisplayshortskip 6\p@ \@plus3\p@ \@minus3\p@
\end{lstlisting}
